% ECONOMICS_PROBLEM: {"id": "C65-37", "title": "Uniqueness for an oscillatory generator of linear growth", "jel_code": "C65", "source_id": "OP-0564"}
% FIRST_POSTED: 2026-10-10
% ATTEMPT_STATUS: solved
\documentclass[11pt]{article}
\usepackage[margin=1in]{geometry}
\usepackage{amsmath,amssymb,amsthm,lmodern,microtype}
\usepackage[colorlinks=true,linkcolor=blue,citecolor=blue,urlcolor=blue]{hyperref}
\hypersetup{pdftitle={Nonuniqueness for a BSDE with an oscillatory generator},pdfauthor={Ji Ho Bae}}
\newcommand{\E}{\mathbb E}
\newcommand{\Prob}{\mathbb P}
\newcommand{\R}{\mathbb R}
\newtheorem{theorem}{Theorem}
\title{Nonuniqueness for a BSDE with generator\texorpdfstring{\\$|z|\sin|z|$}{ |z| sin |z|}}
\author{Ji Ho Bae}
\date{10 October 2026}
% Generated with OpenAI Codex assistance. The exact serving revision was not exposed.
% Hand proof only; no Lean verification is claimed.
\begin{document}
\maketitle

\begin{abstract}
We construct two solutions with the same terminal value for the scalar
Brownian BSDE with generator $g(z)=|z|\sin|z|$. Both solutions have
integrable supremum powers of every finite order. This gives a negative
answer to Fan, Hu and Tang's Problem~5.3.
\end{abstract}

Fix $T>0$ and the augmented natural filtration of a one-dimensional
standard Brownian motion $B$. Fan, Hu and Tang ask whether
\begin{equation}\label{bsde}
 Y_t=\xi+\int_t^T |Z_s|\sin|Z_s|\,ds-\int_t^T Z_s\,dB_s
\end{equation}
has a unique solution when $\xi\in L^p$, $p>1$, and $(|Y_t|^p)$ is of
class~(D) \cite[Problem~5.3, equation~(5.2)]{FHT}. Here class~(D) means
that the stopped variables $|Y_\tau|^p$, $\tau\le T$, form a uniformly
integrable family.

\begin{theorem}\label{main}
For every $T>0$, there are an $\mathcal F_T$-measurable
$\xi\in\bigcap_{1\le p<\infty}L^p$ and two solutions
$(Y^+,Z^+)$, $(Y^-,Z^-)$ of \eqref{bsde} such that, for every finite
$p\ge1$,
\[
 \E\sup_{t\le T}|Y_t^\pm|^p+
 \E\left(\int_0^T|Z_t^\pm|^2\,dt\right)^{p/2}<\infty.
\]
Moreover, $Y_t^+>Y_t^-$ for every $t<T$, almost surely.
\end{theorem}

\begin{proof}
We use It\^o's formula, stopped martingale identities and the
Burkholder--Davis--Gundy inequality in their usual continuous-martingale
forms; see \cite[Chapters~IV and IX]{RY}.

\paragraph{1. A positive process with an integrable inverse-square clock.}
Let
\begin{equation}\label{bessel}
 dR_t=dB_t+R_t^{-1}\,dt,\qquad R_0=1,
 \qquad A_t=\int_0^tR_s^{-2}\,ds.
\end{equation}
The stopped Picard construction gives a strong solution up to exit from
$(0,\infty)$, since the coefficients are locally Lipschitz there.
For $n\ge2$, let $\sigma_n$ be its first exit from $(1/n,n)$ and put
$\tau_n=T\wedge\sigma_n$. It\^o's formula gives
\[
 d(R_t^{-1})=-R_t^{-2}\,dB_t,
 \qquad d(R_t^2)=2R_t\,dB_t+3\,dt.
\]
The stopped integrands are bounded, so
$\E R_{\tau_n}^{-1}=1$ and
$\E R_{\tau_n}^{2}=1+3\E\tau_n\le1+3T$. Consequently
\[
 \Prob(\sigma_n\le T)\le \frac1n+\frac{1+3T}{n^2}.
\]
Thus $R$ exists throughout $[0,T]$ and stays strictly positive.

Another application of It\^o's formula shows that
\[
 L_t=R_t^{-1/2}\exp(A_t/8),\qquad
 dL_t=-\frac{L_t}{2R_t}\,dB_t.
\]
In particular $\E L_{\tau_n}=1$. Cauchy--Schwarz and the preceding
second-moment bound yield
\[
 \E e^{A_{\tau_n}/16}
 =\E\bigl[L_{\tau_n}^{1/2}R_{\tau_n}^{1/4}\bigr]
 \le(\E R_{\tau_n}^{1/2})^{1/2}
 \le(1+3T)^{1/8}.
\]
Fatou's lemma therefore gives
\begin{equation}\label{clock}
 \E e^{A_T/16}\le(1+3T)^{1/8}<\infty.
\end{equation}
In particular $A_T$ has moments of every finite order.

\paragraph{2. A nonzero gap that vanishes at the terminal time.}
Let $\varphi$ and $\Phi$ denote the standard normal density and
distribution function. For $0\le h\le T$, $r>0$, define
\begin{equation}\label{u}
 u(h,r)=\frac1{2\sqrt T}\int_0^h\Phi(-r/\sqrt s)\,ds,
 \qquad u(0,r)=0.
\end{equation}
Differentiating for $h,r>0$ gives
\[
 u_h=\frac1{2\sqrt T}\Phi(-r/\sqrt h),\qquad
 u_{rr}=\frac1{\sqrt T}\Phi(-r/\sqrt h),\qquad
 u_r=-\frac1{2\sqrt T}\int_0^h
             \frac{\varphi(r/\sqrt s)}{\sqrt s}\,ds.
\]
For the second identity, integrate
$\partial_s\Phi(-r/\sqrt s)
 =r\varphi(r/\sqrt s)/(2s^{3/2})$.
These differentiations are justified by exponential decay as $s\downarrow0$.
Thus $u_h=\tfrac12u_{rr}$, $-1\le u_r\le0$, and
$0\le u\le\sqrt T/4$; the derivative bound follows from
$\varphi\le1$ and $h\le T$.

Set
\[
 D_t=u(T-t,R_t),\qquad q_t=u_r(T-t,R_t),
 \qquad q_T=0.
\]
The processes are continuous, $D_T=0$, and $D_t>0$ for $t<T$.
Using \eqref{bessel} and the heat equation gives
\begin{equation}\label{gap}
 dD_t=\frac{q_t}{R_t}\,dt+q_t\,dB_t.
\end{equation}
The identity extends to $T$ by continuity: $q$ is bounded and every
path of $R$ has a positive minimum on $[0,T]$.

\paragraph{3. Realizing the required secant of the generator.}
For $q\in[-1,0]$, write
\[
 s(q)=\begin{cases}\sin(q/2)/(q/2),&q\ne0,\\1,&q=0,\end{cases}
 \qquad s_0=2\sin(1/2).
\]
Then $0<s_0\le s(q)\le1$ and $\cos(q/2)>0$.
For $r>0$ put $k(r)=\lceil(2\pi s_0r)^{-1}\rceil$.
On the interval
\[
 I_r=[2\pi k(r)+\pi/2,\,(2k(r)+1)\pi]
\]
the function
\[
 F(a,q)=\sin a\cos(q/2)+a\cos a\,s(q)
\]
is strictly decreasing, because
\[
 F_a=(\cos(q/2)+s(q))\cos a-a\,s(q)\sin a<0.
\]
Its left endpoint value is $\cos(q/2)>0$, and its right endpoint
value is at most $-1/r$. There is therefore a unique
$a=a(r,q)\in I_r$ with $F(a,q)=-1/r$. This choice is Borel:
on each Borel set where $k(r)$ is constant, the strictly monotone
inverse depends continuously on $(r,q)$. It satisfies
\begin{equation}\label{abound}
 0<a(r,q)\le \frac1{s_0r}+3\pi.
\end{equation}
Both $a+q/2$ and $a-q/2$ are positive. The addition formula for sine
now gives the exact identity
\begin{equation}\label{secant}
 g(a+q/2)-g(a-q/2)=qF(a,q)=-q/r.
\end{equation}
The formula also holds at $q=0$.

Define predictable processes
\[
 Z_t^\pm=a(R_t,q_t)\pm q_t/2,
 \qquad
 Y_t^-=-\int_0^t g(Z_s^-)\,ds+\int_0^t Z_s^-\,dB_s,
 \qquad Y_t^+=Y_t^-+D_t.
\]
By \eqref{abound}, $|Z_t^\pm|\le C(1+R_t^{-1})$ for a fixed
constant $C$. Hence \eqref{clock} implies
\[
 \E\left(\int_0^T|Z_t^\pm|^2\,dt\right)^{p/2}<\infty
 \qquad(1\le p<\infty).
\]
Since $|g(z)|\le|z|$, Cauchy--Schwarz controls the drift integral,
and the Burkholder--Davis--Gundy inequality controls the stochastic
integral. Thus $\E\sup_{t\le T}|Y_t^-|^p<\infty$ for every finite
$p\ge1$. The same holds for $Y^+$ because $D$ is bounded.

Finally, \eqref{gap} and \eqref{secant} give
\[
 dY_t^+=-g(Z_t^-)\,dt+Z_t^-\,dB_t
          +\frac{q_t}{R_t}\,dt+q_t\,dB_t
        =-g(Z_t^+)\,dt+Z_t^+\,dB_t.
\]
Taking $\xi=Y_T^-=Y_T^+$ proves \eqref{bsde} for both pairs.
Their initial values differ by $u(T,1)>0$. For each $p\ge1$, the
integrable variable $\sup_{t\le T}|Y_t^\pm|^p$ dominates all stopped
powers, establishing the required class-(D) condition.
\end{proof}

The same counterexample works for any Brownian dimension $d\ge1$:
use its first coordinate for $B$ and replace each scalar $Z^\pm$ by
$(Z^\pm,0,\ldots,0)$.

\begin{thebibliography}{9}
\bibitem{FHT}
S. Fan, Y. Hu and S. Tang,
\emph{1D nonlinear backward stochastic differential equations:
a unified theory and applications}, arXiv:2309.06233v2,
11 February 2026. Problem~5.3 and equation~(5.2).
\url{https://arxiv.org/html/2309.06233v2#S5}.

\bibitem{RY}
D. Revuz and M. Yor,
\emph{Continuous Martingales and Brownian Motion},
3rd ed., Springer, 1999. Chapters~IV and IX.
\url{https://doi.org/10.1007/978-3-662-06400-9}.
\end{thebibliography}
\section*{Submission and verification}
Author: Ji Ho Bae. Model: GPT Codex; the exact serving revision was not exposed.
The complete hand proof was generated in this session after direct source
checks and an audit of the construction's signs, measurability and moment bounds.
The \href{https://github.com/jbaelaw/oscillatory-bsde-nonuniqueness/blob/071cfc26010f9d3b059f374b8211f3adb628c18c/paper/Proof.pdf}{three-page manuscript},
\href{https://github.com/jbaelaw/oscillatory-bsde-nonuniqueness/blob/071cfc26010f9d3b059f374b8211f3adb628c18c/paper/Proof.tex}{complete original generated source}, and
\href{https://github.com/jbaelaw/oscillatory-bsde-nonuniqueness/blob/071cfc26010f9d3b059f374b8211f3adb628c18c/verification/ProofAudit.md}{proof audit} are retained at a fixed source revision.
Supplementary symbolic and high-precision checks passed; they do not replace
the analytic proof. No Lean formalization or maintainer verification is claimed.

\section{Completion Estimate}
100\%

This is the author's assessment of the stated negative resolution, submitted
for independent expert review. The terminal value and both solutions satisfy
the original power class-(D) condition for every finite $p>1$.

\end{document}
